\documentclass[11pt, a4paper]{article}
\usepackage[utf8]{inputenc}
\usepackage[T1]{fontenc}
\usepackage[spanish]{babel}
\usepackage[margin=2cm]{geometry}
\usepackage{booktabs}
\usepackage{tabularx}
\usepackage{xcolor}
\usepackage{titlesec}
\usepackage{enumitem}
\usepackage{hyperref}
\usepackage{parskip}
\usepackage{fancyhdr}
\usepackage{lmodern}

% Colors
\definecolor{teal}{RGB}{0, 95, 107}
\definecolor{lightgray}{RGB}{245, 245, 245}
\definecolor{darkgray}{RGB}{60, 60, 60}

% Title formatting
\titleformat{\section}{\large\bfseries\color{teal}}{}{0em}{}[\titlerule]
\titleformat{\subsection}{\normalsize\bfseries\color{darkgray}}{}{0em}{}
\titleformat{\subsubsection}{\normalsize\itshape}{}{0em}{}

\titlespacing{\section}{0pt}{12pt}{6pt}
\titlespacing{\subsection}{0pt}{8pt}{4pt}

% Header/footer
\pagestyle{fancy}
\fancyhf{}
\rhead{\small\color{gray}TD7 — Ingeniería de Datos}
\lhead{\small\color{gray}Segundo Parcial}
\rfoot{\small\thepage}
\renewcommand{\headrulewidth}{0.4pt}

% Hyperref config
\hypersetup{
    colorlinks=true,
    linkcolor=teal,
    urlcolor=teal
}

\begin{document}

% Title page
\begin{center}
    {\LARGE\bfseries\color{teal} TD7: Ingeniería de Datos}\\[0.4em]
    {\large Resumen — Segundo Parcial}\\[0.2em]
    {\small Universidad Torcuato Di Tella}
\end{center}

\vspace{0.5em}
\hrule
\vspace{1em}

% =========================================================
\section{8.1 — Calidad y Gobernanza de Datos}

\subsection{Calidad de datos}

\textbf{Definición:} aptitud de un dato para ser usado por quien debe consumirlo.

\textbf{Cómo se mide:}
\begin{itemize}[noitemsep]
    \item \textbf{Análisis univariado:} mín/máx, media, mediana, moda, cuartiles, histogramas, boxplots
    \item \textbf{Análisis bivariado:} coeficiente de correlación, tablas de contingencia, diagramas de dispersión
\end{itemize}

\subsection{8 Dimensiones de calidad}

\begin{tabularx}{\textwidth}{@{}lX@{}}
\toprule
\textbf{Dimensión} & \textbf{Qué mide} \\
\midrule
\textbf{Precisión} & Qué tan cercano a la realidad (ej: ``Perezz'' en vez de ``Perez'') \\
\textbf{Completitud} & No hay NULLs donde no debería haberlos \\
\textbf{Consistencia} & Reglas de negocio se mantienen entre tablas \\
\textbf{Puntualidad} & Disponible cuando se necesita \\
\textbf{Validez} & Tipo correcto, FK válidas, formato correcto \\
\textbf{Unicidad} & No hay duplicados del mismo hecho \\
\textbf{Conformidad} & Formato correcto \\
\textbf{Accesibilidad} & Disponible para los interesados \\
\bottomrule
\end{tabularx}

\subsection{Gobernanza de datos}

Capacidad organizacional de \textbf{asegurar alta calidad} en todos los datos. Implica:
\begin{itemize}[noitemsep]
    \item Definir \textbf{responsables} (\textit{data stewards})
    \item Políticas de almacenamiento, backup y protección
    \item Normas de uso y procedimientos de auditoría
\end{itemize}

\textbf{Ciclo de vida del dato:} Create $\to$ Store $\to$ Use $\to$ Share $\to$ Archive $\to$ Destroy

\subsection{Alcanzando la calidad}

\textbf{Pipelines:} validar datos en cada pasaje entre subsistemas. En streaming, validar continuamente y definir acción ante cada violación.

\textbf{Observabilidad:} comprensión de salud y desempeño de datos. Técnicas:
\begin{itemize}[noitemsep]
    \item Monitoreo automatizado
    \item \textit{Root cause analysis} (RCA)
    \item \textbf{Data lineage:} rastreo de cómo se llega a un dato
\end{itemize}
Objetivo: \textbf{prevenir, detectar y resolver anomalías} lo antes posible.

% =========================================================
\section{8.2 — Privacidad y Compliance}

\subsection{Roles clave}
\begin{itemize}[noitemsep]
    \item \textbf{Data Governance Lead:} fija la estrategia de gobernanza de la organización
    \item \textbf{Data Steward:} ejecuta las tareas diarias de gobernanza
\end{itemize}

\subsection{PII (Personally Identifiable Information)}

Información que \textbf{identifica directamente} a un individuo: nombre, dirección, CUIT, teléfono, e-mail.

Las organizaciones deben cumplir normativas de protección (GDPR, CCPA, etc.). La gobernanza de datos implica rastrear dónde vive y fluye PII en todos los pipelines, y controlar quién tiene acceso.

\subsection{Data Catalogs}

Metadatos curados sobre \textbf{qué} datos tiene la organización, \textbf{quién} los usa y \textbf{dónde} están.

Sirven para:
\begin{itemize}[noitemsep]
    \item Gobernanza y compliance
    \item Rastrear dónde vive y fluye PII
    \item Saber quién tiene acceso a qué datos
    \item Responder: ¿dónde busco mis datos? ¿qué representan? ¿son relevantes?
\end{itemize}

\textbf{Tipos:} in-house (ej: Uber Databook, Airbnb, Netflix), open-source, comerciales.

\subsection{Data Discovery}

Evolución del catálogo: entendimiento \textbf{dinámico y en tiempo real} del estado actual de los datos (no el estado ``ideal catalogado''). Permite democratización del dato — que todos puedan acceder a los datos que necesitan para su trabajo.

\textbf{Niveles de madurez:}
\begin{enumerate}[noitemsep]
    \item \textit{Data reactive:} nadie accede ni usa datos
    \item \textit{Data scaling:} pocos tienen skills para analizar
    \item \textit{Data progressive:} al menos un empleado data-fluent por equipo
    \item \textit{Data fluent:} todos saben cómo acceder a los datos que necesitan
\end{enumerate}

% =========================================================
\section{8.3 — Optimización en SQL Server}

\subsection{Planes de ejecución}

SQL Server genera un plan físico visible en el query optimizer. El flujo de datos va de \textbf{derecha a izquierda} en el diagrama. El porcentaje en cada operador indica su costo relativo.

\subsection{Operadores de acceso a datos}

\begin{tabularx}{\textwidth}{@{}lXX@{}}
\toprule
\textbf{Estructura} & \textbf{Scan} & \textbf{Seek} \\
\midrule
Heap & Table Scan & — \\
Clustered Index & Clustered Index Scan & Clustered Index Seek \\
Nonclustered Index & Index Scan & Index Seek \\
\bottomrule
\end{tabularx}

\vspace{0.5em}
\textbf{Seek} = acceso directo por índice (eficiente). \textbf{Scan} = recorre todo (costoso en tablas grandes).

\subsection{Operadores de agregación}

SQL Server usa dos operadores para \texttt{SUM}, \texttt{AVG}, \texttt{MAX}, \texttt{GROUP BY}, \texttt{DISTINCT}:

\subsubsection{Stream Aggregate}
\begin{itemize}[noitemsep]
    \item \textit{Scalar aggregates} (sin \texttt{GROUP BY}) $\to$ siempre Stream Aggregate
    \item Con \texttt{GROUP BY} $\to$ datos deben venir \textbf{ordenados} (usa índice o inserta Sort operator)
\end{itemize}

\subsubsection{Hash Aggregate (Hash Match)}
\begin{itemize}[noitemsep]
    \item Para tablas grandes con datos \textbf{no ordenados}
    \item No necesita orden; óptimo cuando cardinalidad estimada es pocos grupos
\end{itemize}

\subsubsection{Hash vs Stream}
Si se crea un índice que provee datos ordenados, el optimizador puede elegir Stream en vez de Hash. Si la entrada no está ordenada pero se pide orden explícito, el optimizador puede: insertar Sort + Stream Aggregate, o usar Hash Aggregate + Sort al final.

\subsubsection{Distinct Sort}
\texttt{DISTINCT} puede implementarse con Stream Aggregate, Hash Aggregate o \textbf{Distinct Sort}.

\subsection{Juntas (Joins)}

El optimizador debe decidir: \textbf{orden de las juntas} y \textbf{algoritmo de junta}.

\begin{tabularx}{\textwidth}{@{}lXXllX@{}}
\toprule
\textbf{Algoritmo} & \textbf{Inner table} & \textbf{Outer table} & \textbf{Tamaño} & \textbf{Ordenado} & \textbf{Cláusula} \\
\midrule
\textbf{Hash Join} & Sin índice & Opcional & Cualquiera & No & Equi-join \\
\textbf{Merge Join} & Con índice & Con índice & Grande & \textbf{Sí} & Equi-join \\
\textbf{Nested Loops} & Con índice (must) & Preferable & \textbf{Pequeño} & Opcional & Cualquiera \\
\bottomrule
\end{tabularx}

\vspace{0.5em}
\textbf{Notas clave:}
\begin{itemize}[noitemsep]
    \item \textbf{Nested Loops:} ideal outer pequeño + inner grande indexado. Acepta cualquier predicado de junta.
    \item \textbf{Merge Join:} ambas tablas pre-ordenadas; óptimo con índice clustered o covering en ambas.
    \item \textbf{Hash Join:} no requiere orden; óptimo outer pequeño + inner grande. Implementado con Hash Match.
    \item Merge y Hash \textbf{requieren equi-join}; Nested Loops acepta cualquier predicado.
\end{itemize}

\end{document}
